\section{Proposed Framework}
\label{sec:framework}

The framework introduced in Section~\ref{sec:architecture} is evaluated along
five protocols that together characterise its operating envelope. The
policy-generation protocol (\S\ref{subsec:policy_generation}) describes the
core rule synthesis loop common to every downstream evaluation. Multi-class
attack-family classification (\S\ref{subsec:multiclass}) measures closed-set
accuracy. Zero-day detection via Suspicious Multi-Class Confidence
(\S\ref{subsec:zeroday}) measures graceful degradation on out-of-distribution
traffic. The anonymization ablation (\S\ref{subsec:anon}) measures the
privacy-utility tradeoff under feature-name obfuscation. The adversarial
robustness protocol (\S\ref{subsec:kerckhoffs}) measures evasion cost under a
Kerckhoffs-style threat model. Two LLM rule generators are used throughout:
\texttt{claude-haiku-4-5-20251001} (Anthropic) and \texttt{gemini-2.5-flash}
(Google), enabling a cross-vendor generalisability check on every result.
Global hyperparameters are fixed at $k_{\text{rules}}=5$, $n_{\text{top}}=5$,
$n_{\text{seeds}}=3$, $n_{\text{rounds}}=8$ with seeds $\{42, 123, 456\}$.

\subsection{Policy Generation Protocol}
\label{subsec:policy_generation}

Every reported result derives from a single rule-synthesis loop applied to a
class-balanced corpus per dataset. The loop runs entirely client-side and
emits one policy JSON artefact per (dataset, model, seed) tuple; the artefact
is the unit of reproducibility, and every result in this paper is regenerable
end-to-end from the saved JSON via a single eval invocation.

\paragraph{Class-balanced corpora.}
For each dataset, we sample at most $C = \min(N_{\text{normal}},
N_{\text{attack}}, C_{\max})$ flows per class, where $C_{\max}$ is a
per-dataset cap chosen to bound LLM context size. Bot-IoT inherits the
population's $C = 370$ benign-flow ceiling so the corpus matches the paper's
$740$-sample footprint (Table~\ref{tab:datasets}). UNSW-NB15 uses the
benchmark's pre-split train/test files; all other datasets use a stratified
80/20 split with a 20\% validation slice carved off the training set. The
held-out test set is touched only once at the end of the loop.

\paragraph{Stats package and prompt context.}
Per-feature normal-versus-attack divergence is ranked by the standardised mean
difference, and the top $n_{\text{top}}=15$ features form the
context that the LLM consumes. The human prompt presents (a) the global
attack-versus-normal mean/std/p10/median/p90 for each top feature, (b) per
attack class means for the same features when class labels are available
(this is the cross-class generalisation signal), and (c) value-count
overviews for low-cardinality categorical columns.

\paragraph{Phenomenon-tagged rule proposals.}
Each round, the LLM is asked to emit exactly $k_{\text{rules}}=5$ rules via
a structured \texttt{propose\_rule} tool. Every proposal must carry (i) a
\texttt{feature}, (ii) an \texttt{operator} from
$\{<, \leq, >, \geq, ==, \neq\}$ with explicit constraints prohibiting
equality on continuous floats, (iii) a numeric or categorical \texttt{value},
(iv) a controlled-vocabulary \texttt{phenomenon\_tag} from
$\{\textsc{Handshake}, \textsc{Asymmetry}, \textsc{Timing}, \textsc{Header},
\textsc{Service}, \textsc{State}, \textsc{Volume}, \textsc{Entropy}\}$, and
(v) a one-sentence \texttt{rationale}. The phenomenon vocabulary is the
mechanism by which the next round is told which protocol-level concepts the
current policy already covers, steering the LLM toward uncovered phenomena
and away from per-feature clones.

\paragraph{Diversity-aware acceptance.}
After scoring every candidate rule's per-rule F1 on the validation slice, the
$k$ accepted rules are filled greedily by a composite score
\[
  S(r \mid A) = \alpha \cdot \mathrm{F1}_{\text{val}}(r)
  + (1 - \alpha) \cdot D(r \mid A),
\]
where $A$ is the already-accepted set and the diversity term decomposes as
\[
  D(r \mid A) = w_t \cdot \mathrm{tag}(r,A)
  + w_d \cdot \mathrm{dis}(r,A)
  + w_f \cdot \mathrm{feat}(r,A)
  + w_\theta \cdot \mathrm{thr}(r,A) .
\]
\textit{tag} is one minus the fraction of $A$ already carrying $r$'s
phenomenon tag; \textit{feat} indicates feature novelty; \textit{thr}
measures threshold separation when the feature is re-used; and \textit{dis}
is an imbalance-aware decision disagreement,
$1 - (J_{\text{normal}} + J_{\text{attack}})/2$, with $J_{\text{class}}$ the
Jaccard similarity of fire-sets computed \emph{within} each class. The
per-class decomposition prevents the metric from collapsing to triviality on
severely imbalanced datasets where any two attack-predicting rules agree on
nearly all rows. Default weights are $\alpha = 0.6$ and
$(w_t, w_d, w_f, w_\theta) = (0.35, 0.30, 0.20, 0.15)$.

\paragraph{Weighted voting with calibrated threshold.}
At inference time the policy applies each rule independently to a flow and
combines the binary outputs through a weighted vote
$s(\mathbf{x}) = \sum_i w_i \cdot \mathbb{1}[r_i(\mathbf{x})]$, classifying
$\mathbf{x}$ as attack iff $s(\mathbf{x}) \geq \tau$. Per-rule weights are
the per-rule attack-class precision on the validation slice; $\tau$ is
auto-calibrated on the same slice to maximise the selected metric
(macro-averaged F1 by default; alternatives include attack-class F1 and
attack precision at recall $\geq 0.9$). Both $\{w_i\}$ and $\tau$ are
serialised inside the policy JSON, so the canonical evaluator reproduces any
metric exactly without re-running the LLM.

\paragraph{Round selection and early stopping.}
The loop runs for up to $n_{\text{rounds}} = 8$ rounds and early-stops when
the validation selected-metric does not improve for two consecutive rounds.
The round selected for reporting is the one with the highest
validation-slice score; the test split is never observed during selection.

\subsection{Multi-Class Classification Protocol}
\label{subsec:multiclass}

Multi-class classification is realised through three coupled components. Each
dataset is labelled with its full taxonomy of attack families. Labels are
stored alongside feature vectors in the Chroma vector store so that retrieved
neighbours carry class provenance. The classification prompt lists the
complete set of $K$ known attack classes and instructs the rule generator to
assign one of these labels. Majority voting across the top-$k$ retrieved
neighbours' labels supplies a soft prior to the model as in-context evidence.
Following classification, the rule-generation module synthesises one
decision-rule vector per detected class. Each vector consists of
$k_{\text{rules}}=5$ threshold conditions over the features ranked highest by
permutation importance ($n_{\text{top}}=5$), emitted as JSON tuples
\texttt{(feature\_name, value, op)} for downstream deployment as ACL-style
firewall entries.

\subsection{Zero-Day Detection Protocol}
\label{subsec:zeroday}

Rule-based detectors are closed-set by construction and cannot name a truly
novel attack class. This protocol assesses whether retrieval confidence
degrades gracefully when out-of-distribution traffic appears, enabling an
anomaly flag without a matching rule. For each dataset and a designated
withheld class $c^{*}$, all flows of $c^{*}$ are removed from the training
split and from the Chroma store. The remaining $K-1$ classes form the
closed-set training regime, and $c^{*}$ flows are reintroduced at inference
time alongside balanced known-class samples. Rather than asking the model to
label unseen traffic, we monitor the maximum softmax confidence of the
retrieved neighbour distribution. For a query $\mathbf{x}$, let $p_j$ denote
the proportion of retrieved neighbours labelled with class $j$. The SMCC
score is
\[
  \mathrm{SMCC}(\mathbf{x}) = \max_{j \in \mathcal{K}} p_j ,
\]
where $\mathcal{K}$ is the set of known classes. When
$\mathrm{SMCC}(\mathbf{x}) < \tau$, the query is flagged as a candidate
zero-day event. The threshold $\tau$ is swept from $0$ to $1$ and the
Zero-Day Detection Rate (ZDR) and False Positive Rate are recorded at each
point, producing the curves reported in \S\ref{subsec:threshold}. ZDR is the
fraction of withheld-class samples correctly flagged,
\[
  \mathrm{ZDR} = \frac{\mathrm{TP}_{c^{*}}}{\left| \mathcal{D}_{c^{*}} \right|} ,
\]
where $\mathcal{D}_{c^{*}}$ is the set of test flows belonging to $c^{*}$.

\subsection{Anonymization Ablation}
\label{subsec:anon}

The prompt necessarily exposes raw feature names such as
\texttt{Header\_Length} or \texttt{SrcAddr}, which can leak deployment topology
or vendor information. This ablation tests whether replacing semantic names
with opaque identifiers degrades detection. For each dataset, a bijective map
$\sigma: \text{feature\_name} \to f_i$ is constructed once and held fixed
across seeds. The model receives only the anonymized identifiers and the
numeric values. Two conditions are compared across $n_{\text{seeds}}=3$ random
seeds and $n_{\text{rounds}}=5$ refinement rounds. The \emph{Original}
condition uses semantic names. The \emph{Anonymized} condition uses opaque
identifiers. Performance is reported as macro-averaged F1 on a held-out test
set, and the utility cost of anonymization is
$\Delta F1 = \overline{F1}_{\text{anon}} - \overline{F1}_{\text{orig}}$.
Feature-set consistency across seeds is quantified by the Jaccard index over
the selected $k_{\text{rules}}=5$ features,
\[
  J = \frac{\left| \mathcal{F}_{\text{orig}} \cap \mathcal{F}_{\text{anon}} \right|}
           {\left| \mathcal{F}_{\text{orig}} \cup \mathcal{F}_{\text{anon}} \right|} .
\]

\subsection{Adversarial Robustness Protocol}
\label{subsec:kerckhoffs}

The robustness protocol assumes the Kerckhoffs principle. The attacker has
full knowledge of the generated rule set and applies the minimum feature
perturbation needed to evade detection. The perturbation budget $\varepsilon$
is expressed as a fraction of each feature's interquartile range, and the
adversary perturbs withheld-class flows toward the nearest known-class
centroid. The Exploitability Sweep Rate measures the fraction of
originally-flagged flows that evade detection at budget $\varepsilon$,
\[
  \mathrm{ESR}(\varepsilon) = \frac{\text{flows that evade at budget } \varepsilon}
                                   {\text{flows flagged at } \varepsilon = 0} .
\]
The summary metric is the critical budget $\varepsilon_{0.5}$ at which
$\mathrm{ESR}=0.5$, which quantifies the adversarial hardness of each
detector. Five conditions are compared. \texttt{LLM\_orig} uses semantic
feature names. \texttt{LLM\_anon} uses opaque identifiers. \texttt{DT\_d3}
and \texttt{DT\_d5} are Decision Trees at depths three and five.
\texttt{RF\_d5} is a depth-five Random Forest. Each condition is run for
three seeds, and the mean and standard deviation of $\varepsilon_{0.5}$ are
reported in Table~\ref{tab:kerckhoffs}.
